\begin{theorem}[Two-body form of the components along every generator]
\label{thm:c11-classes}
Let $x$ be a generator and $O=\int\omega\,X_1\cdots X_n$ a monomial local functional, gauge-variant ones included. Let $J=I_{\mathfrak{psl}(4|4)}-I^{CT}_{\mathfrak{psl}(4|4)}[\varepsilon]-\hbar F_1$ with $\hbar F_1$ a one-loop quantisation through arity one along $x$, that is, \eqref{eq:c8-qme} modulo $\hbar^2$ in arity zero and in arity one along $x$, equivalently the first line of \eqref{eq:c9-arity2}, and let $\psi=0$. The endpoint is $T=\infty$ along $\Sigma$, $\Sigma^*$, $\Delta$, $P_b$ and the rotations, and any $T\in(0,\infty)$ along $K_b$.
\begin{itemize}
\item (i) The $\delta$-linear graphs of arity one along $x$ and of order $\hbar$ of the flow of the defect vertex of Lemma \ref{lem:c8-transfer}, applied to $e^{(J+\delta O(\varepsilon))/\hbar}$, form the classes: (F) the one-loop graphs of the vertex $\delta\{S^{(1)}(x),O\}$ and of the vertex $\delta(Q+\{I,-\})O$ with a separate vertex $S^{(1)}(x)$; (C2) the defect cycles of the pairs of letters of $O$, with trees of $I$ on the outputs; (C0) the cycles through two letters of $O$ with the special edge or a contraction through $K_\varepsilon$, with $S^{(1)}(x)$ on a tree; (Br) the graphs with the special edge on a tree; (Ac) the graphs whose cycle passes through vertices of the action alone; (Ct) the trees with one counterterm vertex of order $\hbar$.
\item (ii) (F) is the one-loop part of the arity-one flow of $sO$, and (C0) with the arity-zero counterterm vertices of (Ct) tends to the trees of $\hbar L_1(O)^{(0)}$ with one vertex $S^{(1)}(x)$, both on the left side of \eqref{eq:c8-rep}. A graph of (Br) separates at its special edge, of generic momentum $p$, a one-loop subgraph $G$; when $G$ is divergent the bare graph has the finite limit $\sum_m\frac{(-1)^m}{m!}|p|^{2m}a_m(G)$, $a_m(G)$ the coefficients of $\varepsilon^{-m}$ of $G$, and its counterterm partner, the special-edge part of the counterterm bracket across the same edge, has the opposite limit: the sum is $(e^{-\varepsilon|p|^2}-1)(G-\mathrm{Sing}\,G)$ and tends to $0$. The partners are the special-edge parts of $-\{I,\mathrm{ct}^{(1)}_x(O)[\varepsilon]\}_\varepsilon$, $-\{S^{(1)}(x),\partial_\delta I^{CT}[\varepsilon](O)\}_\varepsilon$ and $-\{I^{CT}[\varepsilon],O\}_\varepsilon$, whose coefficients of $\varepsilon^0$ are the finite special vertices of Definition \ref{def:c11-defect} (iii) and the vacuum-pole terms of \eqref{eq:c10-finitesp}; they are used up here. (Ac) with the defect vertices of the action tends to $0$ by the master equation in arity zero and in arity one along $x$.
\item (iii) (Master formula.) With the Koszul signs of \eqref{eq:c10-decomp},
\begin{equation}
\begin{aligned}
\hbar L_1(O)^{(1)}(x)=\ &\sum_{1\le a<b\le n}(-1)^{|X_b|(|X_{a+1}|+\cdots+|X_{b-1}|)}\int\omega\,X_1\cdots X_{a-1}\,\Upsilon_x(X_a,X_b)\,X_{a+1}\cdots\widehat{X_b}\cdots X_n\\
&-\big\{\hbar F_1^{(1)}(x),O\big\},
\end{aligned}
\label{eq:c11-decompx}
\end{equation}
the two-body kernels $\Upsilon_x=[\varepsilon^0]\mathcal D_x$ of Definition \ref{def:c11-defect} summed over the unordered pairs of letters, and a one-body term, the bracket with the arity-one finite counterterm along $x$. In arity zero, with the same signs,
\begin{equation}
\hbar L_1(O)^{(0)}=\sum_{a<b}\pm\int\omega\cdots\Upsilon_0(X_a,X_b)\cdots-\big\{\hbar F_1^{(0)},O\big\}.
\label{eq:c11-decomp0}
\end{equation}
Every number of a component is a coefficient of $\varepsilon^0$ of a defect cycle or a coefficient of the action-side solution $\hbar F_1$; these formulas are the definition of the components in the construction.
\item (iv) The coefficients of $\log\varepsilon$ and of $\varepsilon^{-m}$ cancel. The defect cycles carry no $\log\varepsilon$ (Theorem \ref{thm:c11-oddsimplex}), and the first cancellation is the logarithmic covariance
\begin{equation}
\mathrm{lc}\big(\{S^{(1)}(x),O\}\big)-\big\{S^{(1)}(x),\mathrm{lc}(O)\big\}=s\,\mathrm{lc}^{(1)}_x(O)-\mathrm{lc}^{(1)}_x(sO)+\big\{L^{(1)}(x),O\big\},
\label{eq:c11-logcov}
\end{equation}
with $L^{(1)}(x)$ the coefficient of $\log\varepsilon$ in $I^{CT}_{\mathfrak{psl}(4|4)}[\varepsilon]^{(1)}(x)$, which has no term $\varepsilon^{-m}$ along $\Sigma$, $\Sigma^*$, $\Delta$, $P_b$ and $K_b$. The second is the pole identity, for $m\ge1$,
\begin{equation}
\sum_{a<b}\pm[\varepsilon^{-m}]\,\mathcal D_x(X_a,X_b)=s\,\mathrm{ct}^{(1)}_{x,m}(O)+\big\{S^{(1)}(x),\mathrm{ct}_m(O)\big\}-\mathrm{ct}_m\big(\{S^{(1)}(x),O\}\big)-\mathrm{ct}^{(1)}_{x,m}(sO),
\label{eq:c11-pole}
\end{equation}
the special-edge parts of the counterterm brackets, with all their powers of $\varepsilon$, having been grouped with the bridges in (ii); for a single pair $O=\int\omega\,X_aX_b$ the right side contains the poles of the pairs of the cubic monomials of $\{S^{(1)}(x),O\}$ and of $sO$ times their third letter.
\item (v) \eqref{eq:c11-decompx} holds along $K_b$ at every $T\in(0,\infty)$ with the same right side. Given the master equation through arity one along $x$, the local representative of $L(O)$ exists through arity one along $x$ for every local functional $O$: (iii) of Definition \ref{def:c10-hyp} holds along every generator.
\end{itemize}
\end{theorem}

\begin{proof}
(i) At order $\hbar$ a graph of the flow has one cycle and vertices of order $\hbar^0$, or no cycle and one vertex of order $\hbar$ (Proposition \ref{prop:c3-graphs}); the part of arity one along $x$ and linear in $\delta$ places $x$ on one vertex and $\delta$ on one vertex. Expand the defect vertex in $\delta$, in the arity and in $\hbar$, with $I_{\mathfrak{psl}(4|4)}=I+S^{(1)}+S^{(2)}$ and $O(\varepsilon)=O-\partial_\delta I^{CT}_{\mathfrak{psl}(4|4)}[\varepsilon](O)$. The classical master equation \eqref{eq:c2-mc} leaves at order $\hbar^0$ the vertices $\delta\{S^{(1)}(x),O\}$ and $\delta(Q+\{I,-\})O$ and the special-edge parts of $\{S^{(1)}(x),O\}_\varepsilon$, $\{I,O\}_\varepsilon$, $\{I,S^{(1)}(x)\}_\varepsilon$, $\tfrac12\{I,I\}_\varepsilon$, together with $\delta\hbar\Delta_\varepsilon O$; $\hbar\Delta_\varepsilon S^{(1)}(x)$ vanishes, its quadratic terms pairing a field with the antifield of another field and its cubic terms carrying $\lambda^*$ and no $\lambda$. A cycle through a vertex of the first two kinds is in (F). A cycle with a special edge passes through $O$ along two letters (Theorem \ref{thm:c10-decomp} (i)) or avoids $O$, and the positions of $S^{(1)}(x)$ and of the special edge sort it into (C2), (C0), (Br) or (Ac).

(ii) By Theorem \ref{thm:c8-L} (iv) the left side of \eqref{eq:c8-rep} in arity one along $x$ is $\big(L(O)^{(1)}(x)\big)[T]+\big(L(O)^{(0)}\big)[T]^{(1)}(x)$ with $L(O)=sO+\hbar L_1(O)$; the one-loop part of $(sO)[T]^{(1)}(x)$ consists of the graphs (F) and their counterterm trees. For (C0), the arity-zero analysis of Theorem \ref{thm:c10-decomp} and of the present argument at arity zero applies with a tree carrying $S^{(1)}(x)$ attached at generic momenta. (Br): the special edge on a tree edge of momentum $p$ multiplies the subgraph $G$ it separates by $e^{-\varepsilon|p|^2}-1=\sum_{j\ge1}(-\varepsilon|p|^2)^j/j!$. For $G=\sum_{m\ge1}\varepsilon^{-m}a_m+a_0+o(1)$, a cycle through two letters of $O$ with $S^{(1)}(x)$ on it or a cycle of the action such as the self-energy with the pole of $I^{CT}$, the terms $j=m$ give the finite limit stated. The counterterm vertex of $G$ joined by the same special edge is $-(e^{-\varepsilon|p|^2}-1)\sum_m\varepsilon^{-m}a_m$; its coefficient of $\varepsilon^0$ is $-\sum_m\frac{(-1)^m}{m!}|p|^{2m}a_m$, which in the bracket notation of \eqref{eq:c10-finitesp} is $-\frac{(-1)^m}{m!}\{f,g_m\}_{H^m}$ with $(f,g_m)=(I,\mathrm{ct}^{(1)}_{x,m}(O))$, $(S^{(1)}(x),\mathrm{ct}_m(O))$ or $(\text{coefficient of }\varepsilon^{-m}\text{ in }I^{CT}[\varepsilon],O)$. The sum of a bridge and its partner is $(e^{-\varepsilon|p|^2}-1)$ times the renormalised subgraph, which has a limit (Definition \ref{def:c8-quant}), and tends to $0$; the partners with $\log\varepsilon$ carry $(e^{-\varepsilon|p|^2}-1)\log\varepsilon\to0$ (Theorem \ref{thm:c8-obstruction} (ii), Proposition \ref{prop:c4-defect1}). Along $K_b$ at finite $T$ the subgraphs carry at fixed $\varepsilon$ the term $\tfrac3{8\pi^2}\hbar\log T$ times a tree with the same factor (remark after Theorem \ref{thm:c8-K}), which also tends to $0$. (Ac): the cycle and its counterterm vertices form the flow of the defect of the action in arity zero or in arity one along $x$, joined to $O$ by a tree at generic momenta, whose limit is the tree-level flow of $\hbar\,\mathrm{Ob}-s(\hbar F_1)$ in these arities (Definition \ref{def:c8-ob}), zero by the master equation.

(iii) What remains is (C2) with the arity-one vertices of (Ct) without special edge, against the tree of $\hbar L_1(O)^{(1)}(x)$ and the counterterm trees of $-\partial_\delta I^{CT}[\varepsilon](\{S^{(1)}(x),O\})$ and $-\mathrm{ct}^{(1)}_x(sO)$ on the left. The vertex $-\{\hbar F_1^{(1)}(x),O\}_\varepsilon$ gives the one-body term, its special-edge part tending to $0$. Each defect cycle is a sum of poles, $[\varepsilon^0]\mathcal D_x$ and a vanishing remainder (Theorem \ref{thm:c11-oddsimplex}); the trees on its outputs converge, and Theorem \ref{thm:c8-L} (i) identifies the limit with the tree-level flow of \eqref{eq:c11-decompx}. The sum over the pairs with the Koszul signs is that of Theorem \ref{thm:c10-decomp} (iii), each defect cycle meeting $O$ in exactly two letters. In arity zero the same argument with the special-edge cycles of the pairs, the vertex $-\{\hbar F_1^{(0)},O\}_\varepsilon$ and the bridges through the self-energy cycles grouped with the special-edge parts of $-\{I^{CT}[\varepsilon],O\}_\varepsilon$ gives \eqref{eq:c11-decomp0}.

(iv) The left side has a limit; collecting the coefficients of $\log\varepsilon$ and of $\varepsilon^{-m}$ of (C2), of the brackets without special edge and of the counterterm trees moved from the left side gives \eqref{eq:c11-logcov} and \eqref{eq:c11-pole}, the special-edge parts of the counterterm brackets cancelling against the bridges at every power of $\varepsilon$ by (ii). A coefficient of $\varepsilon^{-m}$ in $I^{CT}_{\mathfrak{psl}(4|4)}[\varepsilon]^{(1)}(x)$ is a colour-singlet local functional of degree $-1$ and total Weyl weight $4-\tfrac12\lambda_x-2m$, the position at weight $-1$ and of degree $p$ at most that of $x$ (Proposition \ref{prop:c10-bound1}); a colour singlet of degree $-1$ contains an antifield of weight at least $\tfrac52$ and a second adjoint field of weight at least $1$, or $c^*$ and $c$, so its density has weight at least $\tfrac72$ and the total weight is at least $\tfrac72-p$. For $m\ge1$ the total weight is at most $\tfrac52$, $\tfrac32$, $2$, $3$, $1$ along $\Sigma$, $\Sigma^*$, $\Delta$, $P_b$, $K_b$, below $\tfrac72$, $\tfrac52$, $\tfrac52$, $\tfrac72$, $\tfrac32$. Read for one pair, the cubic terms of $\{S^{(1)}(x),O\}$ and of $sO$ have pairs with the empty leg set of Theorem \ref{thm:c10-decomp} (iv), whose poles stand on the right with the third letter.

(v) The argument uses $sO$ through (F), which cancels between the two sides, and through the counterterm trees. Along $K_b$ at a fixed finite $T$ the edges $P(\varepsilon,T)$ are bounded uniformly in $T$ and decay as $e^{-|y|^2/4T}$ (Definition \ref{def:c8-PT}); $\Upsilon_{K_b}$ is independent of $T$ by Theorem \ref{thm:c11-oddsimplex} (iii), and Theorem \ref{thm:c8-K} gives the independence of the local representative. The classes (C2) and (Ct) have local limits, so Theorem \ref{thm:c8-L} (i) gives the local representative.
\end{proof}
